I don't exactly have any ``open questions'' this week, per se, so I'll take the time to
use this section as a space for reflections. It seems as though, when the latent space is
``simple'' enough (e.g. MNIST), the TNBBeta-VAE is comfortable ditching its multimodal
capacity. When the space gets more complicated (i.e. representing graphs with several
connected communities), the model adapts and uses its capacity for expressivity. The
benefit here is that one need not consider whether a unimodal or multimodal posterior is
preferable when making a model selection decision; in the unimodal case, the TNBBeta-VAE
performs just as well as the power-spherical VAE anyway. I see this model as a strict
improvement, at the cost of an MC-estimated KL-divergence.

I'd be curious in exploring other datasets in the wild that need to leverage multimodality.
The fact that two of the SphericalTNBBeta modes are located at opposite poles of the unit
hypersphere gives some signal here. A dataset such as positive/negative sentiment reviews
could be useful in that respect.
